\section{Combinanción y unión de Data Frames}

\subsection{Concatenación \textit{.concat()}}

Esta función es utilizada para unir dos o más \textit{Data Frames} o \textit{Series} a lo largo de un eje, ya sea poniendolo debajo del otro (fila) o al lado del otro (columna).

\subsubsection{Concatenación vertical \texttt{(axis=0)}}

Es el que viene por defecto. En este caso los datos se van a unir de forma vertical, es decir uno abajo del otro.
\begin{pycode}{Unión vertical}
import pandas as pd

df1 = pd.DataFrame({'A': ['A0', 'A1'], 'B': ['B0', 'B1']})
df2 = pd.DataFrame({'A': ['A2', 'A3'], 'B': ['B2', 'B3']})

# Los pegamos uno abajo del otro
resultado = pd.concat([df1, df2])
print(resultado) # Retorna:
#     A   B
# 0  A0  B0
# 1  A1  B1
# 0  A2  B2
# 1  A3  B3
\end{pycode}

\begin{nota}
    Note como los indices se mantienen de forma original. Es decir si \texttt{df1} tiene índices \texttt{0, 1} y \texttt{df2} también tiene \texttt{0, 1}. El resultado tendrá los índices repetidos \texttt{(0, 1, 0, 1)}. Para que Pandas cree un nuevo índice continuo \texttt{(0, 1, 2, 3)}, usa \texttt{ignore\_index=True} de la siguiente forma:
    \begin{verbatim}
        pd.concat([df1, df2], ignore_index=True)
    \end{verbatim}
\end{nota}

\subsubsection{Concatenacion Horizontal \texttt{(axis=1)}}

Cuando queremos unir columnas.
\begin{pycode}{Unión Horizontal}
import pandas as pd

df1 = pd.DataFrame({'A': ['A0', 'A1'], 'B': ['B0', 'B1']})
df2 = pd.DataFrame({'A': ['A2', 'A3'], 'B': ['B2', 'B3']})

# Los pegamos uno abajo del otro
resultado = pd.concat([df1, df2], axis=1)
print(resultado) # Retorna:
#     A   B   A   B
# 0  A0  B0  A2  B2
# 1  A1  B1  A3  B3
\end{pycode}

\subsubsection{¿Qué sucede ante faltantes?}

Si los \textit{Data Frames} no tienen las mismas columnas o filas, por defecto se aplica \texttt{how='outer'}, el cual conserva todas las columnas y en caso de que no exista un dato lo rellena con \texttt{NaN}. (Outer es la \((\cup )\) matemática)

Si usted no quiere conservar datos nulos, puede utilizar \texttt{how='inner'}. (Inner es la intersección \((\cap)\) matemática).

\begin{pycode}{Utilizacion de inner o outer}
import pandas as pd

df1 = pd.DataFrame({'ID': [1, 2], 'Nombre': ['Ana', 'Juan']})
df2 = pd.DataFrame({'ID': [2, 3], 'Sueldo': [500, 600]})

# INNER: Solo queda el ID 2 (el único en ambos)
inner = pd.merge(df1, df2, on='ID', how='inner')
print(inner) # Retorna
#    ID Nombre  Sueldo
# 0   2   Juan     500

# OUTER: Quedan los IDs 1, 2 y 3 (rellena con NaN donde falta)
outer = pd.merge(df1, df2, on='ID', how='outer')
print(outer) # Retorna
#    ID Nombre  Sueldo
# 0   1    Ana     NaN
# 1   2   Juan   500.0
# 2   3    NaN   600.0  
\end{pycode}

\paragraph{Left Join (\texttt{how='left'})}

Este mantiene todas las filas de la tabla izquierda y trae las concidencias de la tambla derecha. Si no hay concidencia rellena los huecos con \texttt{NaN}.
\begin{pycode}{Left join}
import pandas as pd

clientes = pd.DataFrame({'ID': [1, 2, 3], 'Nombre': ['Ana', 'Juan', 'Pedro']})
pedidos = pd.DataFrame({'ID': [1, 2], 'Producto': ['Laptop', 'Mouse']})

# Pedro (ID 3) no tiene pedidos, pero con 'left' se mantiene en la lista
resultado = pd.merge(clientes, pedidos, on='ID', how='left')
print(resultado) # Retorna:
#    ID Nombre Producto
# 0   1    Ana   Laptop
# 1   2   Juan    Mouse
# 2   3  Pedro      NaN
\end{pycode}

\paragraph{Right join (\texttt{how='right'})}

Es lo opuesto al anterior.

\begin{pycode}{Right Join}
import pandas as pd

clientes = pd.DataFrame({'ID': [1, 2, 3], 'Nombre': ['Ana', 'Juan', 'Pedro']})
pedidos = pd.DataFrame({'ID': [1, 2], 'Producto': ['Laptop', 'Mouse']})

# Pedro (ID 3) no tiene pedidos, pero con 'left' se mantiene en la lista
resultado = pd.merge(clientes, pedidos, on='ID', how='right')
print(resultado) # Retorna:
#    ID Nombre Producto
# 0   1    Ana   Laptop
# 1   2   Juan    Mouse
\end{pycode}


\subsubsection{Concatenación con llaves}

Al "pegar" tablas una debajo de otra, a veces perdemos el rastro de qué fila pertenecía a qué archivo original. Para evitar esto, \textit{Pandas} permite usar el parámetro \texttt{keys}, que crea un índice jerárquico o \textit{MultiIndex}.

Esto transforma su tabla en una estructura de niveles: el primer nivel indica el origen (la llave) y el segundo nivel conserva el índice original de la fila.
\begin{pycode}{Concatenación con llaves}
import pandas as pd

df1 = pd.DataFrame({'A': ['A0', 'A1'], 'B': ['B0', 'B1']})
df2 = pd.DataFrame({'A': ['A2', 'A3'], 'B': ['B2', 'B3']})

resultado = pd.concat([df1, df2], keys=['Enero', 'Febrero'])
# Esto te permite hacer: resultado.loc['Enero']

print(resultado) # Retorna:
#             A   B
# Enero   0  A0  B0
#         1  A1  B1
# Febrero 0  A2  B2
#         1  A3  B3
\end{pycode}





\subsection{Unión por \texttt{merge(), join()}}

Ambos métodos sirven para combinar DataFrames basándose en etiquetas o columnas comunes.

\subsubsection{Unión por \texttt{.merge()}}

Esta función une \textit{Data Frames} usando columnas, es ideal cuando tenemos múltiples tablas con columnas comunes. 

Básicamente \texttt{.merge} busca valores iguales en una columna de la Tabla A y una columna de la Tabla B, y cuando los encuentra, pega las filas para crear una sola fila más larga.

\begin{pycode}{.merge()}
import pandas as pd

clientes = pd.DataFrame({'ID': [1, 2, 3], 'Nombre': ['Ana', 'Juan', 'Pedro']})
pedidos = pd.DataFrame({'ID': [1, 2], 'Producto': ['Laptop', 'Mouse']})

# Utilizamos merge
resultado = pd.merge(clientes, pedidos, on='ID')
print(resultado) # Retorna: 
#    ID  Nombre  Producto 
# 0    1    Ana    Laptop 
# 1    2   Juan     Mouse 
\end{pycode}

No siempre tendremos la suerte de que las columnas ``llave'' se llamen igual en ambas tablas. Para estos casos, en lugar de usar el parámetro \texttt{on}, le indicamos explícitamente a \textit{Pandas} qué columna mirar en la tabla de la izquierda (\texttt{left\_on}) y cuál en la de la derecha (\texttt{right\_on}).

\begin{pycode}{Utilización de 'on'}
import pandas as pd

df1 = pd.DataFrame({'A': ['A0', 'A1'], 'B': ['B0', 'B1']})
df2 = pd.DataFrame({'C': ['A1', 'A3'], 'D': ['B1', 'B3']})

resultado = pd.merge(df1, df2, left_on='A', right_on='C')
print(resultado) # Retorna:
#     A   B   C   D
# 0  A1  B1  A1  B1

# Tanto para A como para C el valor es A1    
\end{pycode}




\subsubsection{Unión por (\texttt{.join()})}

El método \texttt{.join()} es una forma rápida y conveniente de combinar dos DataFrames. A diferencia de \texttt{merge()}, su comportamiento está diseñado para unir tablas a través de sus índices (las etiquetas de las filas) en ves de columnas comúnes. Es la version simplificada.

Cuando se utiliza \texttt{df1.join(df2)}, \textit{Pandas} mira los índices de ambos. Cuando encuentra números o etiquetas iguales, pega las columnas de \texttt{df2} al lado de las de \texttt{df1}.

\begin{pycode}{Utilización básica de .join()}
import pandas as pd

# Tabla de empleados con ID como índice
empleados = pd.DataFrame({
    'Nombre': ['Juan', 'Ana']}, 
    index=[1, 2]
)

# Tabla de sueldos con ID como índice
sueldos = pd.DataFrame({
    'Sueldo': [1000, 2000]}, 
    index=[1, 2]
)

# Unimos directamente
resultado = empleados.join(sueldos)
print(resultado) # Retorna:
#   Nombre  Sueldo
# 1   Juan    1000
# 2    Ana    2000 
\end{pycode}

\begin{nota}
    Si ambos \textit{Data Frames} tienen una columna con la misma etiqueta, \texttt{.join()} da un error, al menos que se le diga qué sufijo ponerle a cada uno. Por ejemplo:
    \begin{verbatim}
        df1.join(df2, lsuffix='_izq', rsuffix='_der')
    \end{verbatim}
\end{nota}


En caso de que se quiera unir una columna con un índice, se puede utilizar:
\begin{verbatim}
    df1.join(df2, on='ID_en_df1') 
\end{verbatim}
'on' es el que hace que \texttt{.join()} utilize la columna de \texttt{df1} para buscar el índice de \texttt{df2}.





\subsection{Apilado y Des-apilado (\texttt{stack, unstack})}

\subsubsection{Apilado (\texttt{.stack})}

La función \texttt{.stack()} pasa las columnas a filas. Se dice que ``comprime'' el \textit{Data Frame} para volverlos más largo y estrecho.

\begin{pycode}{Apilamiento}
import pandas as pd

datos = {
    'Notas':[8, 9],
    'Edad':[15, 16]
}

df = pd.DataFrame(datos, index=['Juan', 'Ana'])
print(df) # Retorna:
#       Notas  Edad
# Juan      8    15
# Ana       9    16

print(df.stack()) # Retorna:
# Juan  Notas     8
#       Edad     15
# Ana   Notas     9
#       Edad     16    
\end{pycode}

En el ejemplo, el \textit{Data Frame} tiene 2 indices filas. 


\subsubsection{Desapilar (\texttt{.unstack()})}

Es la operación inversa de \texttt{unstack()}. Este toma un nivel del índice de las filas y lo mueve hacia las columnas.

\begin{pycode}{Desapilado}
import pandas as pd

datos = {
    'Notas':[8, 9],
    'Edad':[15, 16]
}

df = pd.DataFrame(datos, index=['Juan', 'Ana']).stack()

print(df) # Retorna:
# Juan  Notas     8
#       Edad     15
# Ana   Notas     9
#       Edad     16

print(df.unstack()) # Retorna:
#       Notas  Edad
# Juan      8    15
# Ana       9    16    
\end{pycode}
